\subsection{Late normalized sections and the extra open-subgroup datum}

The constructions of \S49.V--VI continue beyond C3 scan~650:
scan~651 starts folio~628 and scans~689--690 are folio~647
and its verso.  In this subsection groups are discrete or
profinite, with continuous maps in the latter case.
We do not identify their automorphism functors with integral
group schemes.  The finite examples may also be read as
constant finite group schemes.  Their constant \(C_2\) is
not replaced by \(\mu_2\) in characteristic two.

\begin{statement}{B.17}{Refuted}{The smaller antidiagonal quotient needs an intersection}
The definition in \S49.VI, fols.~631 verso--632 verso,
C3:658 (18), 659 (28), and 660 (29)--(32), cannot be used
literally with all the data it permits.  Its loop subgroup
\(L_0^\natural\) need not lie in the given open normal
\(C^\natural\).  In that case its antidiagonal does not
lie in \(C^\natural\rtimes Z\), so the displayed smaller
quotient is undefined.  This occurs already in the modular
finite quotient \(C=\operatorname{SL}_2(\mathbb Z/4)\),
with \(C^\natural\) the image of \(\Pi_0\).

Here is the exact repair.  Let \(C\) be a group, or a
profinite group, and let \(D\) be a group of automorphisms
with retained numerical coordinates.  In the profinite case
assume \(D\) is profinite, its action is continuous, and the
map \(j\) below is continuous.  Let \(A\triangleleft C\)
and \(L\subset C\) be \(D\)-stable subgroups, closed in the
profinite case.  Suppose
\[
 j:L\longrightarrow D,\qquad
 j(l)=(\operatorname{Int}(l),1,1,1)
\]
is defined, and \(D\)'s conjugation on \(j(L)\) agrees with
its given action on \(L\).  Write
\[
 \Delta L=\{(l^{-1},j(l)):l\in L\},\quad
 G=(C\rtimes D)/\Delta L,\quad
 H_A=\operatorname{im}(A\rtimes D\longrightarrow G),\quad
 K=L\cap Z(C).
\]
Then the image construction is well-defined and
\[
 H_A\simeq(A\rtimes D)/\Delta(L\cap A),\qquad
 1\longrightarrow A/(A\cap K)\longrightarrow H_A
 \longrightarrow D/j(L\cap A)\longrightarrow1.       \tag{B.17.1}
\]
If \(K=1\), the kernel of \(H_A\to D/j(L)\) is
the image of \(AL\), and it equals the image of \(A\)
if and only if \(L\subset A\).  Thus replacing the
quotient by its image does not also preserve the manuscript's
claimed smaller exact sequence without this extra condition.

Suppose more specifically that
\(D\subset\operatorname{Aut}(C)\times T\times\{\pm1\}^2\),
where \(T\) is the multiplier group.  In the profinite
case take \(T\) profinite and this embedding continuous
and closed.  Put
\[
 P_A=\{(f,m)\in A\times T:
             (\operatorname{Int}(f),m,1,1)\in D\}.
\]
For the homomorphism forgetting the multiplier but retaining
the automorphism and both signs, the exact kernel is
\[
 \ker\bigl(H_A\longrightarrow
       \operatorname{Aut}(C)\times\{\pm1\}^2\bigr)
 \simeq P_A/\{(l,1):l\in L\cap A\}.                 \tag{B.17.2}
\]
For the manuscript's specified normalizer-with-multiplier
\(P_A=S\widetilde N_0^\natural\), this gives the correctly
typed form of C3:664 (52 bis).  Its denominator can be all
of \(L_0^\natural\) only when that subgroup is contained
in \(C^\natural\).
\end{statement}

\begin{proof}
We use the multiplication
\((c,d)(c',d')=(c\,d(c'),dd')\).
The map \(l\mapsto(l^{-1},j(l))\) is a homomorphism:
its first coordinate on a product is \((ll')^{-1}\).
Its image centralizes the copy of \(C\), and conjugation
by \(D\) carries its value at \(l\) to its value at
\(d(l)\).  Hence \(\Delta L\) is normal.
Its intersection with \(A\rtimes D\) is exactly
\(\Delta(L\cap A)\), proving the first isomorphism.
This intersection is normal in the smaller semidirect product.
Projection onto \(D/j(L\cap A)\) gives the exact sequence;
the kernel of \(A\to H_A\) is \(A\cap\ker j=A\cap K\).
When \(K=1\), a representative \((a,d)\) maps trivially
to \(D/j(L)\) precisely when \(d=j(l)\); its image in
\(G\) is then the image of \(al\).  This proves the
assertion about \(AL\).  All subgroups and quotients are
closed profinite groups when the data are profinite, since
they are continuous images or closed subgroups of compact
groups.  No quotient-topology identification is omitted.

For the second assertion use the injective homomorphism
\[
 P_A\longrightarrow A\rtimes D,
 \qquad(f,m)\longmapsto
 (f^{-1},(\operatorname{Int}(f),m,1,1)).
\]
Multiplying two displayed pairs gives the value at
\((ff',mm')\).  Its image is precisely the inverse image
of the asserted automorphism-and-sign kernel.  Its
intersection with \(\Delta L\) consists of
\((l,1)\) with \(l\in L\cap A\), proving (B.17.2).

For the finite modular witness take, modulo~4,
\[
 \rho=\begin{pmatrix}0&1\\-1&1\end{pmatrix},\quad
 \sigma=\begin{pmatrix}0&-1\\1&0\end{pmatrix},\quad
 e_0=\sigma\rho,\quad l=e_0^2=
 \begin{pmatrix}1&2\\0&1\end{pmatrix},
\]
and
\[
 A=\left\{I,
 \begin{pmatrix}1&2\\2&1\end{pmatrix},
 \begin{pmatrix}3&0\\2&3\end{pmatrix},
 \begin{pmatrix}3&2\\0&3\end{pmatrix}\right\}.
\]
This is the normal Klein four subgroup used in E.15--E.16.
Here \(l\notin A\), but \(h=-l\in A\), and
\(\operatorname{Int}(l)=\operatorname{Int}(h)\).
Consequently the tuple \((\operatorname{Int}(l),1,1,1)\)
satisfies C3:653 (4): both conjugators can be \(h\), the
loop is fixed, \(A\) is stable, and the induced action
on \(C/A\) is already trivial.  It therefore satisfies
both the stronger and the weakened versions of condition~d).
The involution \(\tau=\operatorname{diag}(-1,1)\)
normalizes \(A\), as required by the orientation extension.

The same obstruction occurs in
\(C=\widehat{\operatorname{SL}_2(\mathbb Z)}\) itself.
Take \(A=\widehat\Pi_0\), \(l=e_0^2\), and
\(h=-e_0^2\).  By E.15, \(A\) is normal of index~12
and \(h\) is one of its specified generators; the
modulo-four calculation proves \(l\notin A\).
Conjugation by \(\tau\) inverts the two generators of
\(\Pi_0\), so it preserves its closure.  Again
\(\operatorname{Int}(l)=\operatorname{Int}(h)\),
and both conjugators may be \(h\).  Thus
\(l\in L_A\setminus A\) for the actual named modular
subgroup, not only for an abstract finite model.

The program \texttt{B\_late\_sections.py} enumerates all
48 matrices, verifies generation by \(\rho,\sigma\),
normality and these identities, and checks the actual
positive-sign loop preimage in the four-element
\(\langle e_0\rangle\).  It is \(L=\{1,l\}\).
In checking the weakened condition~d), the image of \(\Pi\)
is the eight-element kernel of reduction modulo~2.
Moreover \(L\cap A=L\cap Z(C)=1\), so the example does
not exploit a failure of the earlier injectivity condition.
For the numerical subgroup
\(D=\langle(\operatorname{Int}(l),1,1,1)\rangle\),
the corrected image is \(AL\), of order~8, whereas
\(|A|=4\) and \(D/j(L)=1\).  Its automorphism kernel
is \(\{I,-I\}\).  These are executed checks of both
the quotient obstruction and the exact repair.
\end{proof}

\begin{statement}{B.18}{Open}{The late normalized sections and their remaining profinite comparisons}
The following dictionary records the additional constructions
of \S49.V--VI.  It does not infer the existence of the actual
profinite normalizer sections from the finite cocycle calculation
in D.8.  B.16 supplies normalized conjugators and quotient
sections in its precisely defined \(G_Z\); the further
unquotiented sections and source identifications below require
additional arguments.  We distinguish the smaller fibre product of scan~648,
written \(\mathcal B\) below, from the later open-subgroup
construction depending on \(A=C^\natural\).
The superscript \(\natural\mid\sharp\) is a working reading,
not a palaeographical conclusion.
\end{statement}

\begin{longtable}{@{}p{.19\linewidth}p{.57\linewidth}p{.17\linewidth}@{}}
\toprule Object or map&Definition or assertion to be compared&Source: C3 scan; folio\\\midrule\endhead
\(a_0,b\)&For normalized conjugators,
\(a_0=\alpha^{-1}u\), \(b=\beta^{-1}u\).
The \(b\)-coordinate is multiplicative; the image of \(a_0\)
has the twisted law (47), not an ordinary product law.
&583--585 (41)--(48); 595--596\\
\(A(u),S_0\Gamma^!\)&The fibre over \(u\) consists of
\(a=f'^{-1}u,b=f^{-1}u\), with
\(f'=\alpha\rho^j(\tau')^{\nu_3}\),
\(f=\beta\sigma^i\), \(j\bmod6,i\bmod4\).
The normalized section chooses \(i=j=0\); now
\(a=(\tau')^{\nu_3}\alpha^{-1}u\).
&586--587 (49)--(53); 596v--597\\
\(S_0^\Pi\Gamma^!\)&The congruence-defined subgroup in (62)
is asserted to extend \(Z^!\) by the first central \(C_2\).
The normalized section has cocycle (63); the formula for
\(\alpha''\) requires the final commutator correction.
&590--595 (59)--(64); 598v--601\\
\(M^0(-\bar j),M'''(-\bar j)\)&
\(M^0(-\bar j)=\ker(\epsilon_2|M'(-\bar j))\);
\(M'''(-\bar j)=\langle\tau'\rangle\,M^0(-\bar j)\).
The subgroup lattice, normality, and its stated indices
require the actual normalized sections.
&635--642 (32)--(44); 621--623v\\
\(\Theta,\ker\Theta\)&The map
\(N_\rho^*\to\operatorname{GL}_2(\mathbb Z/4)/V_4\)
is asserted to have kernel \(M^0(-\bar j)\) and quotient
of order~24.  Computing the target is not a proof of this kernel.
&643--647 (45)--(46); 624--626\\
\(E_\rho,S_\sigma,\mathcal B\)&
\(E_\rho=N_\rho^?\to\gamma\) has kernel \(\mu_\rho\);
\(S_\sigma=Z_{\sigma0}^!\to\gamma\) is an isomorphism;
\(\mathcal B=E_\rho\times_\gamma S_\sigma\times_\gamma G\).
These two normalizer maps are substantive inputs.
&647--648 (50)--(55); 626--626v\\
\(\Gamma^\flat\), bisection&
\(\Gamma^\flat=E_\rho\times_\gamma S_\sigma\subset\Gamma\);
its kernel over \(\gamma\) is \(\mu_\rho\).
The resulting canonical map lands in
\(\mathcal M_{\rho,\sigma}/\mu_\rho\).
&649--650 (56)--(60); 627--627v\\
\(A,C^\natural_\tau,Z^A\)&An open normal, \(\tau\)-stable
\(A\subset C\) is extra data; conjugators in (4) belong
to \(A\rtimes\langle\tau\rangle\), with their stated
orientation signs.  The surviving condition~d) requires
stability of \(A\) and identity on
\((C/A)/\operatorname{Im}\Pi\).
&651--653 (1)--(4); 628--629\\
\(L_A,\gamma_A,H_A\)&
\(L_A=\{l\in L_0:j(l)\in Z^A\}\),
\(\gamma_A=Z^A/j(L_A)\).
The image construction \(H_A\) and its smaller kernel
use B.17, not an undefined antidiagonal quotient.
&658--660 (18)--(32); 631v--632v\\
\(S\widetilde N_0^A\)&Pairs \((f,m)\) with
\(f\in A\), \(\operatorname{Int}(f)(e_0)=e_0^m\),
and \((\operatorname{Int}(f),m,1,1)\in Z^A\).
Forgetting the multiplier gives the kernel in (B.17.2).
&663--664 (45)--(52 bis); 634--634v\\
\(M^A,G^A\)&In (107), the three images of
\(\rho,\sigma,e_0\) are \(A\)-conjugate to their specified
\(\tau\)-twists or powers, and \(u\) stabilizes \(A\).
The proposed identifications (108), (115)--(116) require
matching these conditions with the repaired construction.
&682--686 (107)--(118); 643v--645v\\
\(\mathcal M^{A,?}\)&Intersection of the later
\(\mathcal M^A\) with \(\mathcal M^?\).  Its image in
\(M^A\) is determined by the images of
\(N_\rho^{A,?}\) and \(Z_{\sigma0}^{A,!}\) in
\(\Gamma_0^A\); the claimed surjectivity is not checked
in the manuscript.
&688--690 (122)--(128); 646v--647v\\
\bottomrule
\end{longtable}

\begin{proof}[Proved reductions, executed finite cases, and exact residuals]
\emph{The normalized multiplication laws.}
Write \(t=\tau_\gamma\), \(t^2=1\), and suppose
\(\delta(u)=t^{\nu(u)}\delta(a_0(u))\), with
\(\nu(u'u)=\nu(u')+\nu(u)\).  Multiplication alone gives
\[
 \delta(a_0(u'u))=
 t^{\nu(u)}\delta(a_0(u'))t^{-\nu(u)}\delta(a_0(u)).
                                                               \tag{B.18.1}
\]
Thus the untwisted law is equivalent to the commutation
condition in scan~585 (48).  In the finite group \(S_3\),
take \(t=(12)\), \(\nu\) the permutation sign in
\(\mathbb F_2\), and \(a_0(g)=t^{\nu(g)}g\).
The program checks (B.18.1) for all 36 ordered pairs and
exhibits failures of ordinary multiplicativity.

For scan~595 use \(t=\tau'\),
\(a=t^{\nu_3}\alpha^{-1}u\), and the analogous primed
values.  If \(a''=\varepsilon a'a\), with
\(\varepsilon\) central of order dividing two, direct
substitution gives the exact identity
\[
 \alpha''=\varepsilon\,u'(\alpha)\alpha'\,
 t^{\nu_3'}[a',t^{\nu_3}]t^{-\nu_3'}.
                                                               \tag{B.18.2}
\]
Indeed the left side is
\(\varepsilon u'u a^{-1}a'^{-1}t^{\nu_3+\nu_3'}\);
multiplying the right side and cancelling the middle factors
gives that same word.  The program verifies the word identity
for all \(u,u',a,a'\in S_3\) and both sign bits, and finds
a failure when the commutator is omitted.  This is a test
of the stated algebraic identity, not a realization claim
for arbitrary such tuples in the universal profinite group.

The source-order cocycle in scan~595 (63), with the left
factor primed, is
\[
 c((u',v'),(u,v))=u(u'+v')\in\mathbb F_2,
 \quad (\epsilon_3,\epsilon_2)=((-1)^u,(-1)^v).
                                                               \tag{B.18.3}
\]
The associated group on \(\mathbb F_2^3\) has law
\((s,u,v)(t,x,y)=(s+t+x(u+v),u+x,v+y)\).
Its bilinearity proves associativity.  The coordinate change
\((s,u,v)\mapsto(s+uv,u,v)\) identifies it with D.8's
convention.  Hence that classification applies, but its
cocycle must not be copied with the factor order reversed.
The program verifies all 512 associativity triples, the
coordinate change, and all 16 sign pairs.  It confirms the
criterion \(\varepsilon=1\) exactly when
\(\epsilon_3(u)=1\) or
\(\epsilon_2(u')\epsilon_3(u')=1\), thereby correcting
the incompatible prose at the end of scan~594.

\emph{The early bisection, conditional on its typed inputs.}
Let \(N_r,N_s,G\) map onto \(\gamma\), with kernels
\(C_r,C_s,C\).  Suppose subgroups \(E\subset N_r\)
and \(S\subset N_s\), closed in the profinite case, satisfy
\[
 1\to\mu_r\to E\to\gamma\to1,\qquad S\xrightarrow{\sim}\gamma,
                                                               \tag{B.18.4}
\]
where \(\mu_r\subset C_r\) is central in \(N_r\).
Define \(\mathcal M=N_r\times_\gamma N_s\times_\gamma G\),
\(\mathcal B=E\times_\gamma S\times_\gamma G\),
\(\Gamma=N_r\times_\gamma N_s\), and
\(\Gamma^\flat=E\times_\gamma S\).
Pullback of (B.18.4) proves
\[
 1\to\mu_r\to\mathcal B\to G\to1,\quad
 \mathcal B/\mu_r\simeq G\hookrightarrow\mathcal M/\mu_r,
 \quad \Gamma^\flat\cap(C_r\times C_s)=\mu_r\times1.
                                                               \tag{B.18.5}
\]
Every element of \(\Gamma\) differs from an element of
\(\Gamma^\flat\) by an element of its kernel
\(C_r\times C_s\), so
\(\Gamma=\Gamma^\flat(C_r\times C_s)\).
After quotienting by \(\mu_r\), this is a split extension
with kernel \((C_r/\mu_r)\times C_s\) and complement
\(\Gamma^\flat/\mu_r\simeq\gamma\); the action is the
conjugation action of that complement.  This proves exactly
the formal bisection assertions of scans~648--650 under
(B.18.4), without manufacturing the two required
profinite normalizer maps.

If \(\mu_s\subset C_s\), \(\mu_G\subset C\) are the
other central factors, then in the displayed coordinate
order \((r,s,G)\)
\[
 \mathcal B\cap(\mu_r\times\mu_s\times\mu_G)
       =\mu_r\times1\times\mu_G.                    \tag{B.18.6}
\]
Thus the \(1\times\mu\times\mu\) printed in the checked
transcription of scan~648 is incompatible with that order
and with its definition of \(\mu_\rho\).
The second inclusion of scan~647 (50) must have target
\(N_\rho^*\), and the lower right object of
scan~649 (57) must be \(\gamma\), for the displayed kernels.
The final arrow of scan~650 has the canonical target
\(\mathcal M/\mu_r\), as in scan~649 (56); a lift to
\(\mathcal M\) needs a splitting not supplied by the
fibre-product argument.

For an executed finite case take \(\gamma=C_2^2\),
the source-order \(D_8\) extension \(E\) above,
\(N_r=E\times C_3\), \(N_s=\gamma\times C_4\),
\(S=\gamma\times\{0\}\), and \(G=\gamma\times C_2\).
Then \(\mathcal M\simeq E\times C_3\times C_4\times C_2\)
has order~192, \(\mathcal B\) has order~16, and its
kernel over the order-eight group \(G\) has order~2.
The program verifies (B.18.5)--(B.18.6), the order-96
group \(\Gamma\), its order-eight bisection and their
order-two intersection.  Both lifts of the element
\(((1,0),0)\in G\) have order~4.  Hence
\(\mathcal B\to G\) has no section: omitting the quotient
in the purported canonical lift would fail even in this
finite fibre-product model.  These groups are a model of
the formal hypotheses, not asserted to be normalizers
in the same universal modular group.

Finally, for any subgroup restrictions \(E_A\subset E\),
\(S_A\subset S\), and \(G_A\subset G\), the image of
the third projection is exactly
\[
 \operatorname{im}(E_A\times_\gamma S_A\times_\gamma G_A
       \longrightarrow G_A)
 =q^{-1}\bigl(\operatorname{im}E_A\cap\operatorname{im}S_A\bigr).
                                                               \tag{B.18.7}
\]
Both inclusions follow by choosing the two lifts of the
same image in \(\gamma\).  This proves the reduction in
scan~690: surjectivity of the two normalizer maps would
suffice, but must itself be established.

\emph{Exact residuals.}
For the universal profinite modular input, B.16 constructs
the normalized conjugators and the maps
\(Z^!\to N_\rho^*/\langle\omega\rangle\) and
\(Z^!\to N_\sigma^*\), together with their quotient
sections over \(\gamma\), inside the formal \(G_Z\).
It remains to match the congruence-defined source subgroup
\(S_0^\Pi\Gamma^!\) of scans~583--595 and identify
its central extension with the computed sign cocycle.
Prove the normality and index assertions of scans~635--642
and the precise kernel
\[
 \ker\bigl(N_\rho^*\xrightarrow{\Theta}
       \operatorname{GL}_2(\mathbb Z/4)/V_4\bigr)
          =M^0(-\bar j).                              \tag{B.18.8}
\]
This modulo-four target is already the first finite quotient
needed by that construction.  E.19 and its program compute
the target of order~24; they do not compute the kernel of
this profinite-source map.  In particular, stability alone
does not prove the congruence \(\tau'(x)\equiv x\pmod{\Pi_0}\)
in scan~641 (42); its action on the quotient must be checked.
For the typed inputs (B.18.4), B.16 isolates the remaining
unquotiented lift \(\gamma\to N_\sigma^*\) and the
identification of the central pullback in \(N_\rho^*\)
with the source's \(N_\rho^?\).  Thus its proved
quotient-section maps from scans~612--616 are used, while
the additional assertions are retained as explicit residuals.

For the later extra datum take first
\[
 C=\widehat{\operatorname{SL}_2(\mathbb Z)},\qquad
 A=\widehat\Pi=\ker(C\to\operatorname{SL}_2(\mathbb F_2)),
\]
the first named proper subgroup in scan~656, of index~6,
and then \(A=\widehat\Pi_0\), of index~12.
Use \(\tau\), not \(\tau'=\sigma^{-1}\tau\), in the
allowed conjugator subgroup.  In particular
\(\tau(\rho)=\sigma\rho^{-1}\sigma^{-1}\), so changing
the condition to \(A\)-conjugacy to \(\rho^{-1}\) is
not justified when \(\sigma\notin A\).
With the surviving weakened version of condition~d) fixed,
determine the actual \(L_A\), the corrected image \(H_A\),
and the map to the \(M^A\) of scan~682 (107); prove or
correct scan~686 (115)--(116).  Then determine the images
of \(N_\rho^{A,?}\) and \(Z_{\sigma0}^{A,!}\) in
\(\Gamma_0^A\), and therefore the image in (B.18.7),
as requested in scans~689--690.  These are explicit remaining
comparisons for the two named modular subgroups, not a claim
that the corresponding questions are open in the literature.
The \(A=\widehat\Pi_0\) construction already has the
modulo-four obstruction proved in B.17; the original smaller
quotient cannot be silently used in its place.

The finite calculations and formal reductions prove cases
of these exact statements.  They do not prove that the
normalizer subgroups, their images, and their quotient
maps descend to the integral schemes of \S48.
That additional representability and fppf comparison remains
within the explicit scope limitation of A.9 and B.8.
\end{proof}
